% Konfiguracja dokumentu i przełączniki sekcji

% --- Jednostka i logo (puste pole ukrywa element) ---
\newcommand{\University}{Politechnika Wrocławska}
\newcommand{\Faculty}{Wydział Informatyki i Telekomunikacji}
\newcommand{\LogoFile}{logos/logoPWr.png}

% --- Tytuł i przedmiot ---
\newcommand{\DocType}{Dokumentacja projektowa}
\newcommand{\DocTitle}{Analiza poinfekcyjna pamięci RAM oraz logów systemowych na potrzeby inżynierii detekcji}
\newcommand{\Course}{Cyberbezpieczeństwo}
\newcommand{\Subject}{Informatyka Śledcza}
\newcommand{\ShortSubject}{} % opcjonalny skrót nazwy przedmiotu do stopki (puste = używa \Subject)
\newcommand{\GroupInfo}{Grupa nr 2 / wtorek, godz. 11:15 - 14:00}

% --- Autorzy i prowadzący (kolejne linie rozdzielaj przez \\; puste \Supervisor ukrywa blok) ---
\newcommand{\AuthorsLabel}{Autorzy}
\newcommand{\Authors}{%
    Aleh Shehda, nr albumu 272277\\
    Aliaksei Lapachuk, nr albumu 276711
}

\newcommand{\SupervisorLabel}{Prowadzący}
\newcommand{\Supervisor}{%
    dr inż. Robert Czechowski
}

% --- Miejscowość, data, skrót SHA-256 i metadane PDF ---
\newcommand{\PlaceDate}{Wrocław, \today}
\newcommand{\DocumentHash}{4d765532542ad27ee8332d18c750238d7d4ec295ca117fb7aeca216a0339bdf0}
\newcommand{\PdfAuthor}{Aleh Shehda, Aliaksei Lapachuk}

% --- Przełączniki części opcjonalnych (\...true = pokaż, \...false = ukryj) ---
\newif\ifshowabstract         \showabstractfalse
\newif\ifshowabbreviations    \showabbreviationsfalse
\newif\ifshowbibliography     \showbibliographytrue
\newif\ifshowlistoffigures    \showlistoffigurestrue
\newif\ifshowlistoftables     \showlistoftablestrue
\newif\ifshowlistoflistings   \showlistoflistingsfalse
\newif\ifshowlistofalgorithms \showlistofalgorithmsfalse
\newif\ifshowappendices       \showappendicesfalse

% Wersja do druku (\printversiontrue wyłącza kolorowe odnośniki w PDF)
\newif\ifprintversion         \printversionfalse

% Domyślne pozycjonowanie rysunków \Figure (H lub htbp)
\newcommand{\DefaultFigurePlacement}{H}
